\documentclass{article}

% packages
\usepackage{amsmath, amsthm, thmtools, amsfonts, amssymb, luacode, catchfile, tikzducks, hyperref, ifthen}
\ifcsname c@kobocompile\endcsname
	\usepackage[a5paper, total={1072pt, 1448pt}, margin=10pt, includeheadfoot]{geometry} % set page margins
\else
	\usepackage[a4paper, margin=50pt, includeheadfoot]{geometry}
\fi
\usepackage[shortlabels]{enumitem}
\usepackage[skip=3pt, indent=0pt]{parskip}

% language
\usepackage[bidi=basic, layout=tabular, provide=*]{babel}
\ifcsname c@english\endcsname
	\babelprovide[main, import]{english}
\else
	\babelprovide[main, import]{hebrew}
	\babelprovide{rl}
\fi
%\babelfont{rm}{Libertinus Serif}
\babelfont{rm}[Renderer=Harfbuzz]{Libertinus Serif}
\babelfont{sf}{Libertinus Sans}
\babelfont{tt}{Libertinus Mono}

% style
\AddToHook{cmd/section/before}{\clearpage}	% Add line break before section
\linespread{1.3}
\setcounter{secnumdepth}{0}		% Remove default number tags from sections, this won't do well with theorems
\AtBeginDocument{\setlength{\belowdisplayskip}{3pt}}
\AtBeginDocument{\setlength{\abovedisplayskip}{3pt}}
\graphicspath{ {../images/} }

% operators
\DeclareMathOperator\cis{cis}
\DeclareMathOperator\Sp{Sp}
\DeclareMathOperator\tr{tr}
\DeclareMathOperator\im{Im}
\DeclareMathOperator\re{Re}
\DeclareMathOperator\diag{diag}
\DeclareMathOperator*\lowlim{\underline{lim}}
\DeclareMathOperator*\uplim{\overline{lim}}
\DeclareMathOperator\rng{rng}
\DeclareMathOperator\Sym{Sym}
\DeclareMathOperator\Arg{Arg}
\DeclareMathOperator\Log{Log}
\DeclareMathOperator\dom{dom}
\DeclareMathOperator\supp{Supp}
\DeclareMathOperator\var{Var}
\DeclareMathOperator\cov{Cov}

% commands
%\renewcommand\qedsymbol{\textbf{מש''ל}}
%\renewcommand\qedsymbol{\fbox{\emoji{lizard}}}
\newcommand{\Aa}[0]{\mathcal{A}}
\newcommand{\Bb}[0]{\mathcal{B}}
\newcommand{\CC}[0]{\mathbb{C}}
\newcommand{\Cc}[0]{\mathcal{C}}
\newcommand{\EE}[0]{\mathbb{E}}
\newcommand{\FF}[0]{\mathbb{F}}
\newcommand{\Ff}[0]{\mathcal{F}}
\newcommand{\Ii}[0]{\mathcal{I}}
\newcommand{\Gg}[0]{\mathcal{G}}
\newcommand{\Ll}[0]{\mathcal{L}}
\newcommand{\Mm}[0]{\mathcal{M}}
\newcommand{\NN}[0]{\mathbb{N}}
\newcommand{\Nn}[0]{\mathcal{N}}
\newcommand{\PP}[0]{\mathbb{P}}
\newcommand{\Pp}[0]{\mathcal{P}}
\newcommand{\QQ}[0]{\mathbb{Q}}
\newcommand{\RR}[0]{\mathbb{R}}
\newcommand{\Rr}[0]{\mathcal{R}}
\newcommand{\Ss}[0]{\mathcal{S}}
\newcommand{\TT}[0]{\mathbb{T}}
\newcommand{\Uu}[0]{\mathcal{U}}
\newcommand{\Vv}[0]{\mathcal{V}}
\newcommand{\Ww}[0]{\mathcal{W}}
\newcommand{\ZZ}[0]{\mathbb{Z}}
\newcommand{\acts}[0]{\circlearrowright}
\newcommand{\explain}[2] {
	\begin{flalign*}
		 && \text{#2} && \text{#1}
	\end{flalign*}
}
\newcommand{\maketitleprint}[0]{ \begin{center}
	%\begin{tikzpicture}[scale=3]
	%	\duck[graduate=gray!20!black, tassel=red!70!black]
	%\end{tikzpicture}	
	\includegraphics[width=6cm]{cover}
\end{center}
}

% theorem commands
\newtheoremstyle{c_remark}
	{}	% Space above
	{}	% Space below
	{}% Body font
	{}	% Indent amount
	{\bfseries}	% Theorem head font
	{}	% Punctuation after theorem head
	{.5em}	% Space after theorem head
	{\thmname{#1}\thmnumber{ #2}\thmnote{ \normalfont{\text{(#3)}}}}	% head content
\newtheoremstyle{c_definition}
	{3pt}	% Space above
	{3pt}	% Space below
	{}% Body font
	{}	% Indent amount
	{\bfseries}	% Theorem head font
	{}	% Punctuation after theorem head
	{.5em}	% Space after theorem head
	{\thmname{#1}\thmnumber{ #2}\thmnote{ \normalfont{\text{(#3)}}}}	% head content
\newtheoremstyle{c_plain}
	{3pt}	% Space above
	{3pt}	% Space below
	{\itshape}% Body font
	{}	% Indent amount
	{\bfseries}	% Theorem head font
	{}	% Punctuation after theorem head
	{.5em}	% Space after theorem head
	{\thmname{#1}\thmnumber{ #2}\thmnote{ \text{(#3)}}}	% head content

\ifcsname c@english\endcsname
	\theoremstyle{plain}
	\newtheorem{theorem}{Theorem}[section]
	\newtheorem{lemma}[theorem]{Lemma}
	\newtheorem{proposition}[theorem]{Proposition}
	\newtheorem*{proposition*}{Proposition}
	%\newtheorem{corollary}[theorem]{אין חלופה עברית}

	\theoremstyle{definition}
	\newtheorem{definition}[theorem]{Definition}
	\newtheorem*{definition*}{Definition}
	\newtheorem{example}{Example}[section]
	\newtheorem{exercise}{Exercise}[section]

	\theoremstyle{remark}
	\newtheorem*{remark}{Remark}
	\newtheorem*{solution}{Solution}
	\newtheorem{conclusion}[theorem]{Conclusion}
	\newtheorem{notation}[theorem]{Notation}
\else
	\theoremstyle{c_plain}
	\newtheorem{theorem}{משפט}[section]
	\newtheorem{lemma}[theorem]{למה}
	\newtheorem{proposition}[theorem]{טענה}
	\newtheorem*{proposition*}{טענה}
	%\newtheorem{corollary}[theorem]{אין חלופה עברית}

	\theoremstyle{c_definition}
	\newtheorem{definition}[theorem]{הגדרה}
	\newtheorem*{definition*}{הגדרה}
	\newtheorem{example}{דוגמה}[section]
	\newtheorem{exercise}{תרגיל}[section]

	\theoremstyle{c_remark}
	\newtheorem*{remark}{הערה}
	\newtheorem*{solution}{פתרון}
	\newtheorem{conclusion}[theorem]{מסקנה}
	\newtheorem{notation}[theorem]{סימון}
\fi

% Questions related commands
\newcounter{question}
\setcounter{question}{1}
\newcounter{sub_question}
\setcounter{sub_question}{1}

\ifcsname c@english\endcsname
	\newcommand{\question}[1][0]{
		\ifthenelse{#1 = 0}{}{\setcounter{question}{#1}}
		\section{Question \arabic{question}}
		\addtocounter{question}{1}
		\setcounter{sub_question}{1}
	}

	\newcommand{\subquestion}[1][0]{
		\ifthenelse{#1 = 0}{}{\setcounter{sub_question}{#1}}
		\subsection{Part \alph{sub_question}}
		\addtocounter{sub_question}{1}
	}
\else
	\newcommand{\question}[1][0]{
		\ifthenelse{#1 = 0}{}{\setcounter{question}{#1}}
		\section{שאלה \arabic{question}}
		\addtocounter{question}{1}
		\setcounter{sub_question}{1}
	}

	\newcommand{\subquestion}[1][0]{
		\ifthenelse{#1 = 0}{}{\setcounter{sub_question}{#1}}
		\subsection{סעיף \localecounter{letters.gershayim}{sub_question}}
		\addtocounter{sub_question}{1}
	}
\fi

% import lua and start of document
\directlua{common = require ('../common')}

\GetEnv{AUTHOR}

% headers
\author{\AUTHOR}
\date\today

\title{פתרון מטלה 02 --- תורת הקבוצות (80200)}

\begin{document}
\maketitle
\maketitleprint{}

\Question{}
\Subquestion{}
נוכיח שאם $A$ קבוצה סופית וגם $|A| = n$ אז $|\mathcal{P}(A)| = 2^n$.
\begin{proof}
	נקבע ללא פגיעה בכלליות ההוכחה כי הקבוצה $A = [n]$. \\*
	נגדיר פונקציה $f : \mathcal{P}(A) \to [2^n]$ על־ידי
	\[
		f(X) = \sum_{i = 0}^{n} \begin{cases}
			2^i, & i \in X \\
			0, & i \not\in X
		\end{cases}
	\]
	פונקציה זו מתאימה מספר אי־שלילי לכל קבוצה $X \subset A$ על־ידי שימוש בבסיס בינארי, בדומה לייצוג עשרוני, לכן. \\*
	לכן נניח שהפונקציה חד־חד ערכית. לכל מספר $r \in [2^n]$ נוכל למצוא יצוג בינארי מהצורה $r_0 2^0 + r_1 2^1 + \hdots + r_n 2^n$ כאשר $r_i \in \{0, 1\}$ לכל $i \in [n]$. \\*
	אילו נבנה קבוצה $X = \{i \mid r_i = 1\}$ אז נקבל $f(X) = r$ ולכן הפונקציה גם על. \\*
	לכן נובע $|\mathcal{P}(A)| = |[2^n]| = 2^n$
\end{proof}

\Subquestion{}
נוכיח כי אם $|A| = |B|$ אז גם $|\mathcal{P}(A)| = |\mathcal{P}(B)|$.
\begin{proof}
	אילו הקבוצות $A, B$ סופיות אז הטענה נובעת ישירות מהסעיף הקודם. \\*
	נניח כי $A, B$ אינן סופיות. משוויון העוצמות נובע שקיימת $f : A \to B$ הפיכה. \\*
	נגדיר $g : \mathcal{P}(A) \to \mathcal{P}(B)$ על־ידי
	\[
		g(X) = \{ f(x) \mid x \in A \}, \qquad
		g^{-1}(X) = \{ f^{-1}(x) \mid x \in B \}
	\]
	הראינו פונקציה שמוכיחה כי הפונקציה היא על, ומהחד־חד ערכיות של $f$ נובעת חד־חד ערכיות של $g$, לכן $|\mathcal{P}(A)| = |\mathcal{P}(B)|$.
\end{proof}

\Subquestion{}
נוכיח כי מתקיים השוויון
\[
	\mathcal{P}(A) \cap \mathcal{P}(B) = \mathcal{P}(A \cap B)
\]
\begin{proof}
	\[
		\forall X \subseteq A \cap B
		\iff X \subseteq A \land X \subseteq B
		\iff X \in \mathcal{P}(A) \land X \in \mathcal{P}(B)
		\iff X \in \mathcal{P}(A \cap B)
	\]
	ולכן $\mathcal{P}(A) \cap \mathcal{P}(B) = \mathcal{P}(A \cap B)$.
\end{proof}

\Subquestion{}
נוכיח כי
\[
	\mathcal{P}(A) \cup \mathcal{P}(B) \subseteq \mathcal{P}(A \cup B)
\]
\begin{proof}
	\[
		\forall X \in \mathcal{P}(A) \cup \mathcal{P}(B)
		\iff X \subseteq A \lor X \subseteq B
		\implies X \subseteq A \cup B
		\iff X \in \mathcal{P}(A \cup B)
	\]
\end{proof}
נראה דוגמה שבה זו הכלה ממש:
\[
	A = \{0\}, B = \{1\}, 
	\mathcal{P}(A) \cup \mathcal{P}(B) = \{\emptyset, \{0\}, \{1\}\} \subset \mathcal{P}(A \cup B)
	= \{\emptyset, \{0\}, \{1\}, \{0, 1\} \}
\]

\Subquestion{}
נמצא קבוצה אינסופית $A$ כך ש־$A \subseteq \mathcal{P}(A)$. \\*
נגדיר איבר $\emptyset_0 = \emptyset$ ולכל $n \in \NN$ נגדיר $\emptyset_{n + 1} = \{\emptyset_n\}$. \\*
עתה נגדיר $A = \{\emptyset_i \mid i \in \NN \}$. \\*
נובע כי לכל $n$ גם $\emptyset_n \in A \land \emptyset_n \subseteq A$ ולכן גם $\emptyset_n \in \mathcal{P}(A)$ ובמסתכם $A \subseteq \mathcal{P}(A)$.

\Question{}
נבדוק את הרפלקסיביות, הסימטריה והטרנזיטיביות של היחסים הבאים:

\Subquestion{}
\[
	R_1 = \{ \langle x, y \rangle \in \QQ^2 \mid x - y \in \ZZ \}
\]
היחס רפלקסיבי שכן לכל $\langle x, x \rangle$ מתקיים $x - x = 0 \in \ZZ$. \\*
היחס סימטרי שכן $\forall \langle x, y \rangle \in R_1 : x - y \in \ZZ \iff y - x \in \ZZ$. \\*
היחס גם טרנזיטיבי שכן $\forall x, y, z \in R_1 : x R_1 y \land y R_1 z \implies x - y + y - z \in \ZZ, x - z \in \ZZ \implies \langle x, z \rangle \in R_1$.

\Subquestion{}
\[
	R_2 = \{ \langle A, B \rangle \in {(\mathcal{P}(X))}^2 \mid A \cap B \text{ אינסופי} \}
\]
היחס רפלקסיבי שכן $A \cap A = A$ לכל קבוצה, וסימטרי כנביעה ישירה מסימטריית חיתוך הקבוצות. \\*
היחס לא טרנזיטיבי. נגדיר $X = \NN, A = 2\NN, B = \NN, C = 2\NN + 1$ אז $A \cap B = A$ ולכן אינסופית, וגם $B \cap C = C$ אינסופית, אבל $A \cap C = \emptyset$.

\Subquestion{}
\[
	R_2 = \{ \langle A, B \rangle \in {(\mathcal{P}(X))}^2 \mid A \triangle B \text{ אינסופי} \}
\]
היחס לא רפלקסיבי, לדוגמה $\NN \triangle \NN = \emptyset$ והיא איננה אינסופית. \\*
היחס סימטרי כנביעה מסימטריית ההפרש הסימטרי. \\*
היחס גם לא טרנזיטיבי, נראה כי אם $A = C = \NN, B = \emptyset$ אז $A \triangle B = C \triangle B = \NN$ אינסופי אבל $A \triangle C = \emptyset$.

\Question{}
\Subquestion{}
תהי $A$ קבוצה ו־$E \subseteq A \times A$ יחס שקילות. \\*
נוכיח ש־$A / E$ היא חלוקה של $A$.
\begin{proof}
	יהיו $a, b \in A$ ונגדיר $A_a, A_b$ מחלקות השקילות של $a, b$ בהתאמה. \\*
	נשים לב ש־$A_a = A_b \iff aEb$, ולכן נניח מעתה כי $\langle a, b \rangle \not\in E$. \\*
	נקבל בהתאם כי $A_a \ne A_b$, ונראה שגם $A_a \cap A_b = \emptyset$, שאם לא כן קיים איבר $c \in A$ שמקיים $\langle a, c \rangle, \langle b, c \rangle \in E$ ויחד עם טרנזיטיביות $A$ נקבל סתירה. \\*
	כל $a \in A$ מגדיר איזושהי מחלקת שקילות, דהינו אין איבר כך שהוא לא נמצא באף מחלקת שקילות, ולכן איחודן הוא $A$ עצמה. \\*
	מצאנו כי $A/E$ היא חלוקה של $A$.
\end{proof}

\Subquestion{}
תהי $A$ קבוצה ו־$Q \subseteq \mathcal{P}(A)$ חלוקה של $A$, נגדיר
\[
	E_Q = \{ \langle a, b \rangle \in A^2 \mid \exists X \in Q : a \in X \land b \in X \}
\]
נוכיח כי $E_Q$ הוא יחס שקילות.
\begin{proof}
	נבדוק את כל התנאים ליחס שקילות:
	\begin{itemize}
		\item רפלקסיביות: $\forall a \in A \exists X \in Q : a \in X \land a \in X \implies \langle a, a \rangle \in E_Q$.
		\item סימטריה: $\forall \langle a, b \rangle \in E_Q \exists X : b \in X \land a \in X \implies \langle b, a \rangle \in E_Q$.
		\item טרנזיטיביות: $\forall \langle a, b \rangle, \langle b, c \rangle \exists X_1, X_2 : a \in X_1, b \in X_2, b \in X_2, c \in X_2 \implies X_1 = X_2 \implies c \in X_1, \langle a, c \rangle \in E_Q$.
	\end{itemize}
\end{proof}

\Subquestion{}
\subsubsection{i.}
נוכיח כי לכל יחס שקילות $A \subseteq A \times A$ מתקיים $E_{A/E} = E$.
\begin{proof}
	מצאנו בסעיף א' כי $A/E$ היא חלוקה של $A$, ובסעיף ב' מצאנו כי היחס המושרה על־ידי חלוקה מהווה יחס שקילות בין רכיבי החלוקה, ולכן נובע
	\[
		\forall \langle a, b \rangle \in E \iff \langle a, b \rangle \in A_{A/E}
	\]
	ומצאנו כי הטענה נכונה.
\end{proof}

\subsubsection{ii.}
נוכיח כי לכל חלוקה $Q \subseteq \mathcal{P}(A)$ מתקיים $A/E_Q = Q$.
\begin{proof}
	תהי $Q$ חלוקה, מסעיף א' וב' נקבל כי $E_Q$ הוא יחס שקילות וכי $A/E_Q$ חלוקה, נותר להוכיח שהחלוקות אכן זהות. \\*
	יהי $q_1, q_2 \in A$ ונניח כי עבור מחלקת שקילות $Q_1 \in Q$ מתקיים $q_1, q_2 \in Q_1$, לכן מהסעיפים הקודמים נקבל $\langle q_1, q_2 \rangle \in Q_1$ ובהתאם $q_1, q_2$ באותה מחלקת שקילות ב־$A/E_Q$. \\*
	מצאנו כי החלוקות זהות.
\end{proof}

\Question{}
יהי $R \subseteq A \times A$. נוכיח ש־$R$ יחס שקילות אם ורק אם קיימת קבוצה $W$ ופונקציה $f : A \to W$ כך שמתקיים
\[
	R = \{ \langle a, b \rangle \in A^2 \mid f(a) = f(b) \}
\]
\begin{proof}
	\textbf{כיוון ראשון:}
	נניח כי $R$ הוא יחס שקילות ונגדיר $Q_R$ החלוקה המושרית על־ידי $R$. \\*
	נגדיר $W = Q_R$ ו־$f(a)$ מחזיר את מחלקת השקילות $E \in W$ אשר $a \in E$. \\*
	מהגדרה זו נובע כי אם $a R b$ אז גם $f(a) = f(b) = E$ כפי שרצינו.

	\textbf{כיוון שני:}
	נניח כי קיימת קבוצה $W$ ופונקציה $f$ המקיימות את הטענה. \\*
	נראה כי ההגדרה של יחס שקילות מתקיימת:
	\begin{itemize}
		\item רפלקסיביות: $\forall a \in A : f(a) = f(a) \implies \langle a, a \rangle \in R$.
		\item סימטריה: $\forall a, b \in A, \langle a, b \rangle \in R \implies f(a) = f(b) = f(a) \implies \langle b, a \rangle \in R$.
		\item טרנזיטיביות: $\forall \langle a, b \rangle, \langle b, c \rangle \in R: f(a) = f(b), f(b) = f(c) \implies f(a) = f(c) \implies \langle a, c \rangle \in R$
	\end{itemize}
	ומצאנו כי $R$ אכן יחס שקילות.
\end{proof}

\Question{}
תהינה $A, B, C$ קבוצות כך שמתקיים
\[
	|A| < |B|, \qquad |B| < |C|
\]
נוכיח כי $|A| < |C|$.
\begin{proof}
	תהינה פונקציות חד־חד ערכיות
	\[
		f : A \to B, g : B \to C
	\]
	אלה כמובן קיימות על־פי אי־שוויון העוצמות. \\*
	הוכחנו במטלה 1 שהרכבת פונקציות חד־חד ערכיות היא חד־חד ערכית, דהינו $g \circ f : A \to C$ חד־חד ערכית ומתקיים
	\[
		|A| \le |C|
	\]
	ידוע כי $|B| \ne |C|$ ולכן הפונקציה $g$ איננה על, לכן ישנו $x \in C$ כך ש־$\forall b in B, g(b) \ne x$, ולכן גם $\forall a \in A : g(f(a)) \ne x$, זאת אומרת שההרכבה איננה על, ולכן נסיק $|A| \ne |C|$, ומחיבור הטענות:
	\[
		|A| < |C|
	\]
\end{proof}

\Question{}
יהיו $A_1, \hdots, A_k$ קבוצות בנות־מניה.

\Subquestion{}
נניח שהקבוצות זרות בזוגות, ונוכיח כי $\bigcup_{i = 1}^k A_i$ היא בת־מניה.
\begin{proof}
	בתרגול ראינו כי
	\[
		|\NN| = |\overbrace{\NN \times \cdots \times \NN}^k |
	\]
	ולכן יש פונקציה חד־חד ערכית ועל מהטבעיים ל־$k$־יה הסדורה של מספרים טבעיים $f : \NN^k \to \NN$. \\*
	הקבוצה $U = \bigcup_{i = 1}^k A_i$ היא איחוד זר ולכן לכל איבר $u \in U$ קיים $i$ יחיד כך ש־$u \in A_i$, וכי קיימת פונקציה הפיכה $g_i : A_i \to \NN$. \\*
	נגדיר פונקציה $g : U \to \NN^k$ על־ידי $g(a \in A_i) = g_i(a) \cdot e_i$, כאשר $e_i = \langle \overbrace{0, \hdots, 0}^{i - 1}, 1, 0, \hdots \rangle$. \\*
	נקבל כי $f \circ g : U \to \NN$ היא פונקציה הפיכה ולכן $U$ בת־מנייה.
\end{proof}

\Subquestion{}
נוכיח כי $U$ בת־מניה גם כאשר $A_i$ לא זרות.
\begin{proof}
	נגדיר לכל קבוצה $A_i$ קבוצה חדשה $A_i' = \{ \langle i, a \rangle \mid a \in A_i \}$. \\*
	על־פי ההגדרה החדשה $U' = \bigcup_{i = 1}^k A_i'$ הן קבוצות זרות בזוגות, ולכן הקבוצה $U'$ בת־מניה. \\*
	ניצור פונקציה $f : U \to U'$ על־ידי בחירת האינדקס המינימלי אשר מקיים $\langle i, a \rangle \in A_i'$ והגדרת $f(a) = \langle i, a \rangle$. \\*
	הפונקציה המתקבלת היא חד־חד ערכית ולא על, ונסיק $|U| \le |\NN|$. \\*
	נבחר קבוצה ספציפית $A_j \subseteq U$, זוהי קבוצה בת־מניה, ולכן קיימת פונקציה חד־חד ערכית כלשהי $g : \NN \to A_j$, ולכן היא גם חד־חד ערכית $g : \NN \to U$ ככלל, ונסיק $|\NN| \le U$. \\*
	קיבלנו אם כן ש־$|U| = |\NN|$.
\end{proof}

\end{document}
